% 编辑转录副本；不是独立下载原件。来源：https://arxiv.org/html/1806.09179v5

\paragraph{Notation from Section 1.2.}
Let $\mathbb F$ be a finite field, and let $\chi:\mathbb F\to\mathbb C$ be a nontrivial additive character. Given a function $F:X\to\mathbb F$, its bias is
\[
\operatorname{bias}(T):=\mathbb E_{x\in X}[\chi(F(x))].
\]
In particular, let $T:V^d\to\mathbb F$ be an order-$d$ tensor. The bias of $T$ is always real and in $(0,1]$. To see that, define $T(\cdot,x^2,\ldots,x^d)$ to be the order-$1$ tensor on $x^1$ given a fixing of $x^2,\ldots,x^d$. Then
\begin{align*}
\operatorname{bias}(T)&=\mathbb E_{x^2,\ldots,x^d\in V}\left[\mathbb E_{x^1\in V}[\chi(T(x^1,\ldots,x^d))]\right]\\
&=\Pr_{x^2,\ldots,x^d\in V}[T(\cdot,x^2,\ldots,x^d)\equiv0]. \tag{1}
\end{align*}
This is since for an order-$1$ tensor (namely, a linear form), its bias is $1$ if it is identically zero, and is $0$ otherwise. The analytic rank of $T$ is defined to be
\[
\operatorname{arank}(T):=-\log_{|\mathbb F|}\operatorname{bias}(T).
\]
\begin{theorem}[1.5]
Let $T,S:V^d\to\mathbb F$ be order-$d$ tensors. Then
\[
\operatorname{arank}(T+S)\leq\operatorname{arank}(T)+\operatorname{arank}(S).
\]
\end{theorem}
\paragraph{2. Proof of Theorem 1.5.}
We prove Theorem 1.5 in this section. It suffices to prove that for any two order-$d$ tensors $T,S:V^d\to\mathbb F$ it holds that
\[
\operatorname{bias}(T+S)\geq\operatorname{bias}(T)\operatorname{bias}(S). \tag{3}
\]
We first introduce some notation. Let $\mathbf x=(x^1,\ldots,x^d),\mathbf y=(y^1,\ldots,y^d)\in V^d$. For $I\subseteq[d]$ define $I^c=[d]\setminus I$ and shorthand $\mathbf x^I:=(x^i:i\in I)$. Define $T_I(\mathbf x,\mathbf y)$ as
\[
T_I(\mathbf x,\mathbf y):=T_I(\mathbf x^I,\mathbf y^{I^c})=T(z^1,\ldots,z^d),\qquad
z^i=\begin{cases}x^i&\text{if }i\in I,\\y^i&\text{if }i\notin I.\end{cases}
\]
That is, $T_I$ is the tensor $T$ evaluated over $\mathbf x^I,\mathbf y^{I^c}$. Observe that $T(\mathbf x+\mathbf y)$ decomposes as the sum
\[
T(\mathbf x+\mathbf y)=\sum_{I\subseteq[d]}T_I(\mathbf x^I,\mathbf y^{I^c}). \tag{4}
\]
Express $\operatorname{bias}(T)\operatorname{bias}(S)$ as
\begin{align*}
\operatorname{bias}(T)\operatorname{bias}(S)
&=\operatorname{bias}(T(\mathbf x)+S(\mathbf y))\\
&=\operatorname{bias}(T(\mathbf x)+S(\mathbf x+\mathbf y)).
\end{align*}
Here, we used the fact that the joint distributions of $(\mathbf x,\mathbf y)$ and $(\mathbf x,\mathbf x+\mathbf y)$ are identical. Next, decompose $S(\mathbf x+\mathbf y)$ using Equation 4 as
\begin{align*}
\operatorname{bias}(T)\operatorname{bias}(S)
&=\operatorname{bias}\left(T(\mathbf x)+\sum_{I\subseteq[d]}S_I(\mathbf x^I,\mathbf y^{I^c})\right)\\
&=\operatorname{bias}\left((T+S)(\mathbf x)+\sum_{I\subsetneq[d]}S_I(\mathbf x^I,\mathbf y^{I^c})\right).
\end{align*}
Fix a choice of $\mathbf y=\mathbf b\in V^d$ so that
\[
\operatorname{bias}(T)\operatorname{bias}(S)\leq
\left|\operatorname{bias}\left((T+S)(\mathbf x)+\sum_{I\subsetneq[d]}S_I(\mathbf x^I,\mathbf b^{I^c})\right)\right|.
\]
Observe that $S_I(\mathbf x^I,\mathbf b^{I^c})$ is an order-$|I|$ tensor of the inputs $\mathbf x^I$. The proof of Theorem 1.5 follows from the following lemma, applied to $R_{[d]}(\mathbf x)=(T+S)(\mathbf x)$ and $R_I(\mathbf x^I)=S_I(\mathbf x^I,\mathbf b^{I^c})$ for $I\subsetneq[d]$.
\begin{lemma}[2.1]
For each $I\subseteq[d]$, let $R_I:V^I\to\mathbb F$ be an order-$|I|$ tensor. Consider the function
\[
R(\mathbf x)=\sum_{I\subseteq[d]}R_I(\mathbf x^I).
\]
Then $|\operatorname{bias}(R)|\leq\operatorname{bias}(R_{[d]})$.
\end{lemma}
In order to prove Lemma 2.1, we first prove the following claim.
\begin{proposition}[Claim 2.2]
Let $W_0,\ldots,W_n:\mathbb F^m\to\mathbb F$ be functions. Consider functions $A,B:\mathbb F^n\times\mathbb F^m\to\mathbb F$ defined as follows:
\[
A(x,y)=\sum_{i=1}^n x_iW_i(y),\qquad B(x,y)=A(x,y)+W_0(y).
\]
Then $|\operatorname{bias}(B)|\leq\operatorname{bias}(A)$.
\end{proposition}
\begin{proof}
We have
\begin{align*}
\operatorname{bias}(B)&=\mathbb E_{x,y}[\chi(B(x,y))]\\
&=\mathbb E_y[1_{W_1(y)=\ldots=W_n(y)=0}\cdot\chi(W_0(y))].
\end{align*}
Hence
\[
|\operatorname{bias}(B)|\leq\mathbb E_y[1_{W_1(y)=\ldots=W_n(y)=0}]=\operatorname{bias}(A).
\]
\end{proof}
\begin{proof}[Proof of Lemma 2.1]
Fix $i\in[d]$ and decompose $R(\mathbf x)$ as
\[
R(\mathbf x)=\sum_{I\subseteq[d],i\in I}R_I(\mathbf x^I)+\sum_{I\subseteq[d],i\notin I}R_I(\mathbf x^I).
\]
Setting $x=x^i$ and $y=\mathbf x^{[d]\setminus\{i\}}$, the first sum has the form $\sum x_iW_i(y)$, and the second sum which does not depend on $x$ is $W_0(y)$. Claim 2.2 then gives that
\[
|\operatorname{bias}(R)|\leq\operatorname{bias}\left(\sum_{I\subseteq[d],i\in I}R_I(\mathbf x^I)\right).
\]
Applying this iteratively for $i=1,\ldots,d$ completes the proof.
\end{proof}
